\documentclass[11pt]{letter}
\usepackage[T1]{fontenc}
\usepackage[utf8]{inputenc}
\usepackage[margin=1in,top=0.75in]{geometry}
\usepackage{hyperref}
\renewcommand{\opening}[1]{\thispagestyle{firstpage}%
  {\raggedright\toname \\ \toaddress \par}%
  \vspace{0.5\baselineskip}\noindent#1\par\noindent}
\setlength{\longindentation}{0pt}

\signature{Th\'eo Hermann\\Senseable City Lab\\Massachusetts Institute of Technology\\77 Massachusetts Avenue, Cambridge, MA, USA\\thermann@mit.edu}

\address{Th\'eo Hermann\\Senseable City Lab\\Massachusetts Institute of Technology\\77 Massachusetts Avenue\\Cambridge, MA, USA}

\date{\today}

\begin{document}

\begin{letter}{The Editors\\Nature Cities}

\opening{Dear Editors,}

We are pleased to submit our manuscript entitled ``Mapping Leaf Area Index (LAI) in Urban Forests'' for consideration as an Article in \textit{Nature Cities}.

Urban trees are among the most effective tools for mitigating the Urban Heat Island effect, yet cities lack scalable ways to assess canopy shading capacity at the individual-tree level. Leaf Area Index (LAI)---the ratio of leaf area to ground area---is the biophysical parameter most directly linked to a tree's shading potential and, consequently, its local cooling benefit. Traditional arborist-led inventories are labor-intensive, infrequently updated, and systematically exclude private-property trees, which often constitute the majority of the urban forest.

Our manuscript addresses this gap with two complementary contributions. First, we present an automated framework that integrates airborne LiDAR, four-band aerial imagery (RGB + near-infrared), and deep learning--based crown segmentation to detect individual trees and estimate their LAI through voxelization of the three-dimensional canopy structure. Applied to Amsterdam (25~km$^2$), this pipeline detected 12,350 trees---three times the official municipal registry---capturing the private-property canopy that conventional inventories miss, and producing LAI estimates validated against hemispherical photography.

Second, to extend LAI mapping beyond areas with LiDAR coverage, we train a machine learning model that predicts per-tree LAI from morphological parameters and genus. This prediction step enables deployment in municipalities where only aerial imagery and approximate height data are available, lowering the data barrier for adoption. Both the dataset (Hugging Face) and code (GitHub) are publicly available.

The resulting city-wide LAI maps give planners a spatially explicit basis for targeting shade interventions in heat-vulnerable neighborhoods, bridging urban forestry, climate adaptation, and environmental justice. The methodology transfers to any city with standard aerial survey programs. We believe this work aligns closely with \textit{Nature Cities}' interdisciplinary mission.

This manuscript has not been published and is not under consideration elsewhere. All co-authors---Michiel van Selm, Titus Venverloo, Fabio Duarte, and Carlo Ratti---have reviewed and approved submission. We look forward to your consideration.

\closing{Sincerely,}

\end{letter}
\end{document}
